\begin{abstract}
Keystroke signals are attractive triggers for timely help if they change before a learner stops producing text. We
tested whether the 30-second change, at the start of a pause, in two signals of a keystroke state estimator with
hand-set coefficients (the probability of its Stuck state and a smoothed friction axis) distinguishes pauses that became
stalls in character production from matched pauses of 2--30~s that did not. The test was preregistered on a random
holdout of an introductory programming dataset, with an equivalence band of $[0.45, 0.55]$ for the AUC set in advance;
a leak had shown that on seen writers the AUC lay inside it, so this is a confirmation, not a test of a prediction that
could fail. In the holdout (139 students), the AUC fell within the band for both signals (0.489 and 0.487; the
false-equivalence rate is at most 5.2\% [Wilson 95\% interval 3.6--7.5\%] among the configurations examined in the
calibration on seen writers, and is not an upper bound over the whole null). The Stuck-state probability is mostly at
the floor of the estimate, and claims about run-up change rely on the friction axis and on the pairs not at the floor.
In explorations outside the registration, changes of the
observed size barely moved the signals on seen writers, and on the holdout, learners on simple keystroke quantities
ranked stalls higher, mainly through less preceding typing. Under a conversion made after the band was set, a
difference at its edge would change the share of trigger firings that fall on stalls by about one point, so on this
input the signals gave no basis for a trigger. Preregistration: \url{https://osf.io/mrkxt}; results:
\url{https://osf.io/uj6cy}; code: [archive to be added before submission].\\

{\parindent0pt
\textbf{Keywords:} keystroke logging, stalls, equivalence testing, preregistration, matched controls, introductory
programming
}
\end{abstract}
